\documentclass{article}
\usepackage[UTF8]{ctex}
\usepackage{amsmath,amssymb,amsfonts}
\usepackage{enumitem}
\usepackage{geometry}
\geometry{a4paper, margin=1in}
% === 错题本排版优化（不参与题目标识） ===
\setlength{\parindent}{0pt}
\setlength{\parskip}{0.35em plus 0.10em minus 0.08em}
\setlength{\emergencystretch}{3em}
\setlist{itemsep=0.25em, topsep=0.40em, parsep=0pt}
\allowdisplaybreaks[2]
% === 错题本排版优化结束 ===
\begin{document}
\section*{河北工业大学2024年数学分析真题}
\section*{66．河北工业大学2024年数学分析试卷}

\subsection{题目1}
1．求和函数 \(\sum_{n=1}^{\infty} n x^{n-1}=\) \(\_\_\_\_\)．

\subsection{题目2}
2 ．求极限 \(\lim _{x \rightarrow 0} \frac{e^{x} \sin x-x(1+x)}{x^{3}}=\) \(\_\_\_\_\)．

\subsection{题目3}
3．求球面 \(x^{2}+y^{2}+z^{2}=1\) 在点 \((x, y, z)=\left(\frac{\sqrt{2}}{2}, 0, \frac{\sqrt{2}}{2}\right)\) 处的切平面方程 \(\_\_\_\_\)．

\subsection{题目4}
4．求函数 \(f(x)=(x-2)^{\frac{2}{3}}\) 在区间 \([-1,2]\) 上的最大值和最小值 \(\_\_\_\_\)．

\subsection{题目5}
5．设 \(S=\left\{(x, y, z) \mid x^{2}+y^{2}+z^{2}=1\right\}\) ，方向向外，求曲面积分 \(\iint_{S} x \mathrm{~d} y \mathrm{~d} z+y \mathrm{~d} z \mathrm{~d} x+z \mathrm{~d} x \mathrm{~d} y=\)

\subsection{题目6}
6．求极限 \(\lim _{n \rightarrow \infty}\left[\frac{1}{n \ln ^{2} n}+\frac{1}{(n+1) \ln ^{2}(n+1)}+\cdots+\frac{1}{(2 n-1) \ln ^{2}(2 n-1)}\right]=\) 二．计算题．每题 10 分，共 60 分．

\begin{itemize}

\subsection{题目7}
\end{itemize}

\subsection{题目1}
原题边界“1．”：判断函数 \(f(x)=\frac{x^{2}}{4}-\cos x\) 的凸性。

\subsection{题目8}

\subsection{题目2}
原题边界“2．”：设 \(F\) 为二阶可微的三元函数，又设由方程
\[
F(x y, y+z, x z)=0
\]
可以确定隐函数 \(z=f(x, y)\) ，求偏导数 \(\frac{\partial f}{\partial x}\) 和 \(\frac{\partial f}{\partial y}\) ．

\subsection{题目9}

\subsection{题目3}
原题边界“3．”：设曲面 \(\sum\) 为球面 \(x^{2}+y^{2}+z^{2}=1\) 在 \(z^{2} \geq \frac{1}{2}\) 的部分，求 \(I=\iint_{\sum}\left(1-z^{2}\right) \mathrm{d} S\) ．

\subsection{题目10}

\subsection{题目4}
原题边界“4．”：设 \(L\) 是从 \(A\) 点 \((2,0)\) 出发，先沿直线 \(x+2 y=2\) 到 \(B\) 点 \((0,1)\) ，再沿圆周 \(x^{2}+y^{2}=1\) 到 \(C\) 点\((-1,0)\) ，最后由 \(C\) 点沿 \(x\) 轴到 \(A\) 的定向曲线，计算第二类曲线积分
\[
\int_{L}\left(x^{2}-3 y\right) \mathrm{d} x+\left(3 x+e^{y}\right) \mathrm{d} y
\]

\subsection{题目11}

\subsection{题目5}
原题边界“5．”：求解不定积分 \(\int \frac{1}{1+x^{3}} \mathrm{~d} x\) 。

\subsection{题目12}

\subsection{题目6}
原题边界“6．”：计算 \(\int_{0}^{+\infty} \frac{e^{-x^{2}}-e^{-4 x^{2}}}{x^{2}} \mathrm{~d} x\) ．

\subsection{题目13}
1．设 \(f(x)=x^{a}\) ，其中参数 \(0<s<1\) ，证明：\(f(x)\) 在 \([0,+\infty)\) 上一致连续。

\subsection{题目14}
2．设 \(D\) 为平面上一个有界区域，它在 \(x\) 轴与 \(y\) 轴上投影长度分别为 \(a\) 和 \(b, D\) 的面积为 \(A,\left(x_{0}, y_{0}\right) \in\) \(D\) ，证明：

\begin{itemize}
  \item[（1）] \(\left|\iint_{D}\left(x-x_{0}\right)\left(y-y_{0}\right) \mathrm{d} x \mathrm{~d} y\right| \leq a b A\) ．
  \item[（2）] \(\left|\iint_{D}\left(x-x_{0}\right)\left(y-y_{0}\right) \mathrm{d} x \mathrm{~d} y\right| \leq \frac{1}{4} a^{2} b^{2}\) ．
\end{itemize}

\begin{itemize}

\subsection{题目15}
\end{itemize}

\subsection{题目3}
原题边界“3．”：证明函数 \(J(x)=\frac{1}{\pi} \int_{0}^{\pi} \cos (t-x \sin (t)) \mathrm{d} t\) 满足 Bessel 方程
\[
x^{2} J^{\prime \prime}(x)+x J^{\prime}(x)+\left(x^{2}-1\right) J(x)=0 .
\]

\subsection{题目16}

\subsection{题目4}
原题边界“4．”：设函数 \(f(x)\) 在 \([a, b]\) 上连续，证明：\(M(x)=\sup _{a \leq t \leq x} f(t)\) 在 \([a, b]\) 上连续．
\end{document}
